\section{Compact formulation}
\label{sec:model}

This section summarizes N\textsuperscript{2}DBN from an observed event log to
the learned multiresolution graphs, conditional parameters, and generated
event logs. The following sections retain derivations and implementation
details. The construction combines DBN semantics
\citep{dean1989model,murphy2002dbn}, multiscale temporal representations
\citep{saria2007granularity,kang2022multiscale}, coherent probabilistic
allocation \citep{das_topdown2023}, and coupling between related networks
\citep{grzegorczyk2012globalcoupling,kamalabad2020edgewise}.

\paragraph{Event log and temporal granularity.}
\label{sec:model_scales}

An event log is
\(\mathcal D=\{(\ell_i,s_i,d_i)\}_{i=1}^N\), where \(\ell_i\in\mathcal L\)
is an event label, \(s_i\) is its start time, and \(d_i\in\mathbb N\) is its
duration in units of the finest epoch. Let \(k=1,\ldots,K\) index nested
temporal partitions from finest to coarsest, with epochs
\(\mathcal T^{(k)}=\{I_t^{(k)}\}\). Counts and total durations at each
granularity are
\begin{align}
 X_t^{(k,\ell)}
 &=\sum_{i=1}^N\mathbb{I}\{\ell_i=\ell,\ s_i\in I_t^{(k)}\},
 &Y_t^{(k,\ell)}
 &=\sum_{i=1}^Nd_i\mathbb{I}\{\ell_i=\ell,\ s_i\in I_t^{(k)}\}.
 \label{eq:model_log_projection}
\end{align}
Write \(Z^{(k)}=(X^{(k)},Y^{(k)})\). The map \(\pi_k\) assigns an epoch to
its containing epoch at scale \(k+1\). It induces the aggregation map, its
fiber, and event-log reconstruction at the finest granularity,
\begin{align}
 (\mathcal A_kZ^{(k)})_c
 &=\sum_{t:\pi_k(t)=c}Z_t^{(k)},
 &Z^{(k+1)}&=\mathcal A_kZ^{(k)},
 &\mathcal F_k(v)&=\{u:\mathcal A_ku=v\},
 \label{eq:model_aggregation}\\
 \mathcal E(Z^{(1)})
 &=\{(\ell,t,t+Y_t^{(1,\ell)}-1):X_t^{(1,\ell)}=1\}.
 \label{eq:model_episode}
\end{align}
Counts obey the capacity of their epochs, durations obey \(Y\geq X\), and
\(X_t^{(1,\ell)}\in\{0,1\}\).

Let \(\mathcal Q=\{\mathrm{cnt},\mathrm{dur}\}\). A learned edge is
\(e=((a,q_s),(b,q_t),d)\), with \(1\leq d\leq p_k\). For its indicator
\(w_{e,k}\), the lagged values available to target \((b,q)\) are
\begin{align}
 \mathcal P_{k,b,q}(t)
 &=\{Z_{t-d,q'}^{(k,a)}:w_{((a,q'),(b,q),d),k}=1\}.
 \label{eq:model_lag_values}
\end{align}
The learned event structure is
\(\mathcal G=\{\mathcal G^{(k)}\}_{k=1}^K\). Its edges have positive lag, so
every unrolled graph is acyclic. Aggregation and latent links are fixed parts
of the model construction and are excluded from \(\mathcal G\).

Let \(U_t^{(k)}=(S_t^{(k)},H_t^{(k)})\) combine a discrete regimen for
persistent change and a continuous factor for same-epoch co-movement. Context
is denoted by \(C_t^{(k)}\). Resolution-specific maps define the observation
parameters and latent dynamics,
\begin{align}
 \eta_{t,b,q}^{(k)}
 &=f_{k,b,q}(\mathcal P_{k,b,q}(t),U_t^{(k)},C_t^{(k)};\Theta^{(k)}),
 \qquad q\in\mathcal Q,
 \label{eq:model_count_score}\\
 U_t^{(k)}
 &\sim R_k(U_{t-1}^{(k)},U_{\pi_k(t)}^{(k+1)},C_t^{(k)};\Theta^{(k)}).
 \label{eq:model_regime_init}
\end{align}
Boundary arguments are omitted at \(t=1\) and \(k=K\). At the coarsest scale,
an emission map \(E_K\) specifies the count and duration-surplus families. At
finer scales, a normalized allocation map \(W_k\) defines a child kernel,
\begin{align}
 P_K(Z^{(K)}\mid U^{(K)},C^{(K)})
 &=E_K(\boldsymbol\eta^{(K)};\Theta^{(K)}),
 \label{eq:model_coarse_count_mean}\\
 K_k(Z^{(k)}\mid Z^{(k+1)},U^{(k)},C^{(k)})
 &=W_k(\boldsymbol\eta^{(k)},Z^{(k+1)};\Theta^{(k)}),
 \qquad \operatorname{supp}K_k=\mathcal F_k(Z^{(k+1)}).
 \label{eq:model_count_omega}
\end{align}
The implemented \(E_K\) uses zero-inflated negative binomial counts and a
negative binomial duration surplus. The map \(W_k\) uses constrained
Dirichlet-Multinomial allocations for counts and duration surplus. It includes
capacity, sequential depletion, the duration floor, the active-child set, and
the block normalizer.

Let \(\Psi_K\) contain the coarsest latent dynamics and emissions, and let
\(\Psi_k\), \(k<K\), contain the child dynamics and kernel. The joint model is
\begin{align}
 p(\mathcal G,\Theta,Z^{(1:K)},U^{(1:K)}\mid C^{(1:K)})
 &=p(\mathcal G)p(\Theta\mid\mathcal G)\,\Psi_K
   \prod_{k=1}^{K-1}\Psi_k.
 \label{eq:model_joint}
\end{align}
The fiber support gives exact reconciliation for every draw.

\begin{proposition}[Exact multiresolution coherence]
\label{prop:model_reconciliation}
For every parameter value and every draw of the model,
\begin{align}
 \sum_{t:\pi_k(t)=c}X_t^{(k,\ell)}&=X_c^{(k+1,\ell)},
 &\sum_{t:\pi_k(t)=c}Y_t^{(k,\ell)}&=Y_c^{(k+1,\ell)}
 \quad\text{almost surely},
 \label{eq:model_reconciliation}
\end{align}
and consequently \(\mathcal A_{j-1}\cdots\mathcal A_kZ^{(k)}=Z^{(j)}\) for
every \(k<j\).
\end{proposition}

\begin{proof}
\Cref{eq:model_count_omega} gives \(\operatorname{supp}K_k=\mathcal
F_k(Z^{(k+1)})\), and \(u\in\mathcal F_k(v)\) means \(\mathcal A_ku=v\) by
\cref{eq:model_aggregation}, which is the one-step claim. For \(j>k+1\),
suppose \(\mathcal A_{j-2}\cdots\mathcal A_kZ^{(k)}=Z^{(j-1)}\); applying
\(\mathcal A_{j-1}\) and using the one-step case at \(j-1\) gives the claim, so
the identity composes by induction on \(j-k\).
\end{proof}

\Cref{prop:model_reconciliation} is a statement about the arrays \((X,Y)\). The
reconstruction \(\mathcal E\) of \cref{eq:model_episode} additionally imposes a
one-slot floor, a horizon clip and, where labels are declared mutually
exclusive, an occupancy projection. \Cref{sec:identifiability} states each and
bounds the trace-level discrepancy they introduce; the arrays are exact.

\paragraph{What each resolution can identify.}
\label{sec:model_blockconstant}

The kernel \(W_k\) is normalized over a fiber, and that normalization removes a
subspace of its own predictor from the likelihood at scale \(k\). Fix a parent
epoch \(c\in\mathcal T^{(k+1)}\), write \(\child_k(c)=\{t:\pi_k(t)=c\}\), and
decompose the coordinates on that block orthogonally,
\begin{align}
 \sR^{\child_k(c)}
 &=\mathcal N_c^{(k)}\oplus\mathcal V_c^{(k)},
 &\mathcal N_c^{(k)}&=\operatorname{span}\{\mathbf 1_{\child_k(c)}\},
 &\mathcal V_c^{(k)}&=\Bigl\{v:\textstyle\sum_{t\in\child_k(c)}v_t=0\Bigr\},
 \label{eq:model_block_split}
\end{align}
with \(\Pi_k\) the block-averaging projector onto the block-constant
directions,
\((\Pi_kv)_t=|\child_k(\pi_k(t))|^{-1}\sum_{s\in\child_k(\pi_k(t))}v_s\).

\begin{proposition}[Block-constant invariance]
\label{prop:model_block_invariance}
For every parent epoch \(c\), every \(a\in\sR\) and every predictor
\(\boldsymbol\eta^{(k)}\),
\begin{align}
 W_k(\boldsymbol\eta^{(k)}+a\mathbf 1_{\child_k(c)},Z^{(k+1)})
 &=W_k(\boldsymbol\eta^{(k)},Z^{(k+1)}).
 \label{eq:model_block_invariance}
\end{align}
Hence \(\Pi_k\boldsymbol\eta^{(k)}\) is unidentified from \(\Psi_k\), and only
\((I-\Pi_k)\boldsymbol\eta^{(k)}\) is.
\end{proposition}

\begin{proof}
Within block \(c\) the allocation concentrations are
\(\alpha_t=\kappa_k\,\sigma_t(\boldsymbol\eta^{(k)})\) with \(\sigma\) the
softmax over \(\child_k(c)\). The softmax is invariant to a constant added to
every coordinate of the block, so \(\alpha\) is unchanged, and the constrained
Dirichlet-Multinomial density on \(\mathcal F_k(Z^{(k+1)})\) depends on
\(\boldsymbol\eta^{(k)}\) only through \(\alpha\).
\end{proof}

Three consequences organize the model.

First, \(\Pi_k\boldsymbol\eta^{(k)}\) is not lost. It acts on the data through
the parent total \(Z^{(k+1)}\), so it is identified from
\(\Psi_{k+1},\ldots,\Psi_K\). Every parameter has a resolution at which it is
identified, and the two directions of the learning pass carry exactly these two
components: the coarse-to-fine message supplies \(\Pi_k\boldsymbol\eta^{(k)}\)
and the fine-to-coarse message supplies the sufficient statistics for
\((I-\Pi_k)\boldsymbol\eta^{(k)}\). This is the sense in which the pass is
bidirectional.

Second, a term constant across a block contributes exactly zero at scale \(k\).
A latent regimen entering additively is block-constant on every block over
which its path does not switch, so at \(k<K\) it must act on a feature carrying
within-block variation,
\begin{align}
 \eta_{t,b,q}^{(k)}
 &\supseteq\beta_{S_t^{(k)},b,q}^{(k)}\,g_{t,b,q}^{(k)},
 \qquad (I-\Pi_k)g^{(k)}\neq0,
 \label{eq:model_regime_interaction}
\end{align}
with \(g\) prefix-measurable: the depletion feature, a lagged parent value, or
the recency column below. An additive regimen level is admitted at \(k=K\)
alone.

Third, the cross-resolution prior acts on the block-constant component, which
is the component two adjacent scales share.

\paragraph{A covariate that is never block-constant.}

The lag axis is the largest source of competing representations of one
relation, and one prefix-measurable summary of it removes most of that
competition. Let \(\delta_t^{(k,\ell)}=t-\max\{s<t:X_s^{(k,\ell)}>0\}\) be the
recency of label \(\ell\) and \(M_t^{(k,\ell)}\) its geometric-lag counterpart,
\begin{align}
 M_t^{(k,\ell)}(\rho)
 &=(1-\rho)\sum_{d\geq1}\rho^{\,d-1}Z_{t-d}^{(k,\ell)},
 \qquad\text{columns }
 \log\bigl(1+\delta_t^{(k,\ell)}\bigr),\;M_t^{(k,\ell)}.
 \label{eq:model_recency}
\end{align}
Both are measurable with respect to the strict past, so they preserve
acyclicity, and neither is block-constant at an interior scale, so both stay
identified under \cref{prop:model_block_invariance}.

\paragraph{The ladder is a band-pass decomposition.}
\label{sec:model_bands}

Summing over the children of a block is a low-pass filter, so
\cref{prop:model_block_invariance} has a spectral reading. Let \(n_k\) be the
number of finest slots in one scale-\(k\) epoch. A signal supported on
frequencies above \(1/n_{k+1}\) has zero average on every scale-\(k+1\) block
and is therefore pure detail at scale \(k\); a signal supported below
\(1/n_k\) is constant across each scale-\(k\) block and is pure shadow there.
The detail at scale \(k\) is the band
\begin{align}
 \mathcal B_k=\Bigl[\tfrac{1}{n_{k+1}},\ \tfrac{1}{n_k}\Bigr),
 \qquad
 \bigoplus_{k=1}^{K}\mathcal B_k=\text{the whole spectrum},
 \label{eq:model_bands}
\end{align}
so the resolutions partition frequency, and each one estimates the part of the
signal that lives in its own band. Two bases realize this, and the structure decides
which one applies.

\emph{Periodic structure takes a Fourier basis.} A calendar effect of period
\(P\) enters as a few harmonics,
\begin{align}
 \gamma_{P}(t)=\sum_{h=1}^{H}\Bigl[a_h\cos\tfrac{2\pi ht}{P}+b_h\sin\tfrac{2\pi ht}{P}\Bigr],
 \label{eq:model_harmonic}
\end{align}
which costs \(2H\) parameters where a free profile costs \(P-1\), and which is
smooth by construction. Harmonic \(h\) has frequency \(h/P\) and is therefore
identified at the unique scale whose band contains it. \Cref{eq:model_harmonic}
with \(H\) at its maximum recovers the free profile exactly, so the harmonic
order is a smoothness dial and not a modelling commitment. This matters most
where epochs are few: a week-scale effect estimated from nine epochs per
participant cannot afford a free profile and can afford two harmonics.

\emph{Localized structure takes a Haar basis.} A level shift is not periodic,
and a Fourier expansion represents it badly. In the block decomposition of
\cref{eq:model_block_split} the detail coefficients are Haar coefficients on
the block tree, and a level shift at scale \(k\) is exactly a large one,
\begin{align}
 \delta_t^{(k)}=\bigl\langle\eta^{(k)},\psi_t^{(k)}\bigr\rangle,
 \qquad
 \psi_t^{(k)}\in\mathcal V_{\pi_k(t)}^{(k)} .
 \label{eq:model_haar}
\end{align}
So a changepoint has a home scale in the same sense a calendar node does: a
shift that persists longer than one scale-\(k\) block appears as shadow there
and as detail above, and one that decays inside a block appears as detail
there and is invisible above. The latent regimen \(S^{(k)}\) is the
probabilistic version of the same object, with the switch times its
changepoints, and reading them off the posterior at each scale returns the
shifts the ladder can resolve.

\paragraph{Calendar quantities are nodes, and each has a home scale.}
\label{sec:model_calendar}

Hour of day, weekday, day index and week index are vertices of
\(\mathcal G\). Each is a known function of the clock
and has no parents, so an edge from one into a label may carry lag zero without
threatening acyclicity, and the conditional
\(P_{\mathrm{day}}(\text{walk}\mid\text{weekday}=0)\) becomes a table the model
reports. Whether such a node earns its place is decided by the same price as
any other candidate.

\Cref{prop:model_block_invariance} says where each can be earned. A calendar
node \(A\) is a deterministic function of the epoch, so it enters \(\Psi_k\)
only through \((I-\Pi_k)A\); it is therefore identified at scale \(k\) exactly
when it varies inside a scale-\(k\) parent block, and contributes nothing at any
coarser scale where it does not.

\begin{corollary}[Home scale of a calendar node]
\label{cor:model_calendar_home}
Let \(A\) have period \(P_A\) in finest slots and let \(n_k\) be the number of
finest slots in one scale-\(k\) epoch. Then \(A\) is identified at scale \(k\)
if and only if \(n_{k+1}>P_A\) is false and \(n_k<P_A\); equivalently, its home
scale is the finest \(k\) whose parent epoch spans at least one full period of
\(A\).
\end{corollary}

So hour of day is identified where a day is the parent, weekday where a week is
the parent, and a trend index wherever the horizon is the parent. Placing a
calendar node above its home scale buys nothing, and the price it would pay
there is wasted. This turns a modelling convention into a consequence of the
kernel.

For a typed relation \(r\), let \(z_{r,k}\) indicate that some lag is active.
Adjacent structures and standardized active effects evolve according to
\begin{align}
 \Pr(z_{r,k}=1\mid z_{r,k+1}=1)&=\pi_{r,k}^{\mathrm{keep}},
 &\Pr(z_{r,k}=1\mid z_{r,k+1}=0)&=\pi_{r,k}^{\mathrm{appear}},
 \label{eq:model_scale_coupling}\\
 \Pi_k\alpha_{r,k}\mid\alpha_{r,k+1}
 &\sim\mathcal N(\mathcal B_{k+1\to k}(\alpha_{r,k+1}),\tau_{r,k}^2I),
 &(I-\Pi_k)\alpha_{r,k}
 &\sim\mathcal N(0,\tau_{r,k,\perp}^2I).
 \label{eq:model_parameter_evolution}
\end{align}
The tie is on the block-constant component alone, because by
\cref{prop:model_block_invariance} that is the component the parent scale
identifies and the child scale does not. The within-block component has no
parent counterpart and carries its own prior.

The broadcast \(\mathcal B_{k+1\to k}\) is the adjoint of block averaging,
rescaled to the child's feature units,
\begin{align}
 \bigl(\mathcal B_{k+1\to k}\alpha\bigr)_{t}
 &=\frac{\operatorname{sd}(F_{r,k})}{\operatorname{sd}(F_{r,k+1})}\,
 \alpha_{\pi_k(t)}\,\frac{d_{k+1}}{d_k},
 \qquad
 \mathcal B_{k+1\to k}=\Pi_k\mathcal B_{k+1\to k},
 \label{eq:model_broadcast}
\end{align}
where \(d_k\) is the lag at which \(r\) is active at scale \(k\), so the ratio
\(d_{k+1}/d_k\) converts a parent-scale lag into the child's clock. Relations
and lag sets may change by resolution; when \(z_{r,k+1}=0\) the first line of
\cref{eq:model_parameter_evolution} is dropped and \(\alpha_{r,k}\) carries the
second prior alone. A Bernoulli gate on the tie, with a Beta-conjugate update
for \(\pi_{r,k}^{\mathrm{keep}}\) and \(\pi_{r,k}^{\mathrm{appear}}\), lets an
abrupt resolution change escape the coupling; \(\tau_{r,k}\) and the gate are
estimated, and \cref{sec:learning} gives their updates.

\paragraph{One price per relation.}

Scale, lag, quantity pairing and basis are alternative representations of one
scientific relation, and pricing each separately makes a simple relation
compete against its own alternatives. Let \(r=(a,b)\) be an ordered label pair
and let \(g\in\mathcal R(r)\) index a representation of it. The structural
penalty in \cref{eq:model_elbo} is nested,
\begin{align}
 \sum_k\lambda_k|\mathcal E^{(k)}|
 \;\longrightarrow\;
 \sum_k\Bigl[
 \lambda_k^{\mathrm{rel}}\sum_r z_{r,k}
 +\lambda_k^{\mathrm{rep}}\sum_r z_{r,k}\bigl(|\mathcal R_k(r)|-1\bigr)
 \Bigr],
 \qquad
 \lambda_k^{\mathrm{rel}}\gg\lambda_k^{\mathrm{rep}},
 \label{eq:model_group_price}
\end{align}
so a relation is bought once and its further representations are cheap
refinements. The search of \cref{sec:learning} accordingly proposes relations
and reads the lag, quantity and basis out inside the accepted group. Self-lags
\(r=(a,a)\) are a declared part of the construction, exempt from
\(\lambda^{\mathrm{rel}}\) and reported in their own stratum, so a label's own
autocorrelation cannot outbid a cross-label relation for a slot under the
in-degree bound.

Because \(\lambda_k^{\mathrm{rel}}\) prices one purchase out of \(C_k\)
candidate relations, it carries the search's own multiplicity. Under a
marginal-uniform prior on the active set the \(j\)-th accepted relation is
priced at
\begin{align}
 \lambda_k^{\mathrm{eff}}(j)
 &=\lambda_k^{\mathrm{rel}}+\log\frac{C_k-j+1}{j},
 \label{eq:model_multiplicity_price}
\end{align}
which is what makes a per-scale edge count comparable across resolutions whose
candidate fields differ in size.

\paragraph{Bidirectional structural EM.}

For a structured posterior \(q(U^{(1:K)})\), learning maximizes the penalized
evidence lower bound. An interior scale receives evidence from its own factor
and the factor through which it generates its children,
\begin{align}
 \mathcal J(q,\Theta,\mathcal G)
 &=\mathbb E_q[\log p(Z,U,\mathcal G,\Theta\mid C)]
   +\mathbb H[q]
   -\sum_{k=1}^K\lambda_k|\mathcal E^{(k)}|
   -\lambda_\Theta\Omega(\Theta),
 \label{eq:model_elbo}\\
 \log q_k^*(U^{(k)})
 &=\mathbb E_{q_{-k}}[\log\Psi_k+\log\Psi_{k-1}]+\mathrm{constant}.
 \label{eq:model_bidirectional_update}
\end{align}
Fine-to-coarse and coarse-to-fine messages therefore precede the three
structural EM coordinates \citep{friedman1998learning},
\begin{align}
 q^+&=\arg\max_q\mathcal J(q,\Theta,\mathcal G),
 \label{eq:model_em_steps}\\
 (\Theta^{(k)})^+
 &=\arg\max_{\Theta^{(k)}}\left\{
 \mathbb E_{q^+}[\log\Psi_k]+\Lambda_k(\Theta^{(k)})
 -\lambda_\Theta\Omega_k(\Theta^{(k)})\right\},
 \label{eq:model_parameter_step}\\
 \mathcal P_{k,b}^+
 &=\arg\max_{\substack{P\subseteq\mathcal C_{k,b}\\|P|\leq D_k}}
 \left\{\max_{\beta_P}\mathbb E_{q^+}[\log\Psi_{k,b}(P,\beta_P)]
 +\Gamma_k(P)-\lambda_k|P|\right\}.
 \label{eq:model_structure_step}
\end{align}
These are respectively the E-step, M-step, and S-step. Here
\(\Lambda_k(\Theta^{(k)})=\log p(\Theta^{(k)}\mid
\Theta^{(k-1)},\Theta^{(k+1)})\) couples adjacent parameter blocks, and
\(\Gamma_k(P)\) is the adjacent-graph log prior, with unavailable neighbors
omitted at the boundary scales. The S-step profiles edge
coefficients inside each normalized local likelihood. Search may be
greedy or exhaustive under the in-degree bound \(D_k\). Each accepted edge is
checked against shifted versions of its source. Regime count, factor dimension,
and penalty are selected by held-out evidence. A sweep is retained only when
its monitored penalized evidence improves; otherwise the best state is
restored with a smaller M-step learning rate. Cells within one allocation
block share a normalizer and active set, so their exact updates remain
dependent. Participants and conditionally independent blocks can run in
parallel.

Regime permutations, factor-basis changes, and block-constant score shifts
form observational equivalence classes. Parameter recovery is therefore
reported on invariant contrasts and aligned latent states. Missing values are
marginalized within the same fibers. Fitting returns a graph and conditional
parameters at every resolution. Simulation returns one event log whose
projected counts and durations reproduce all generated coarser views.

\paragraph{A population of logs.}
\label{sec:model_participants}

A study observes \(n\) participants, and days inside one participant are
repeated measures. Index participants by \(i\) and let the graph
\(\mathcal G\) and the fixed effects be shared while the level of each
predictor varies,
\begin{align}
 \eta_{i,t,b,q}^{(k)}
 &=f_{k,b,q}\bigl(\mathcal P_{k,b,q}(t),U_{i,t}^{(k)},C_t^{(k)};\Theta^{(k)}\bigr)
 +\bigl(B_i^{(k)}\bigr)_{b,q},
 &B_i^{(k)}&\sim\mathcal N\bigl(0,\Sigma_B^{(k)}\bigr).
 \label{eq:model_participant}
\end{align}
Latent paths are participant-specific; \(\Theta^{(k)}\) and
\(\mathcal G^{(k)}\) are not. Two things follow. A relation shared across
participants borrows strength, which is what makes a per-scale graph estimable
from horizons that are short in epochs. And a label present for some
participants and absent for others cannot enter another label's score through
participant identity, because \(B_i^{(k)}\) absorbs the level that identity
would otherwise supply. Lag histories never cross a participant boundary, and
an epoch whose lagged predictors fall outside the observed interval has its
emission factor masked while the epoch itself is retained, so the sums of
\cref{prop:model_reconciliation} stay intact.

\paragraph{Forecasting.}
\label{sec:model_forecast}

Let \(\mathcal H_T\) be the history through epoch \(T\) at every resolution.
The \(h\)-step predictive distribution is generated by the same factorization,
run forward at the coarsest scale and downward through the child kernels,
\begin{align}
 p\bigl(Z^{(1:K)}_{T+1:T+h}\mid\mathcal H_T,C\bigr)
 &=\int\Bigl[\prod_{s=1}^{h}
 E_K\bigl(\boldsymbol\eta^{(K)}_{T+s}\bigr)
 \prod_{k=K-1}^{1}
 W_k\bigl(\boldsymbol\eta^{(k)}_{T+s},Z^{(k+1)}_{T+s}\bigr)\Bigr]
 \,dq\bigl(U^{(1:K)}_{T+1:T+h}\mid\mathcal H_T\bigr).
 \label{eq:model_forecast}
\end{align}
Latent states are filtered to the origin and propagated by \(R_k\); no
observation after \(T\) is read. Ancestral sampling draws each path by
unrolling \(E_K\) and then descending the chain, and lagged predictors are
replaced by the path's own earlier draws.

\begin{corollary}[Coherent forecasts]
\label{cor:model_coherent_forecast}
Every draw from \cref{eq:model_forecast} satisfies
\(\mathcal A_kZ^{(k)}_{T+s}=Z^{(k+1)}_{T+s}\) for all \(k\) and
\(s\leq h\). The predictive distribution is therefore supported on the coherent
set, and its per-resolution marginals require no reconciliation.
\end{corollary}

\begin{proof}
Immediate from \cref{prop:model_reconciliation}, whose argument uses only the
support of \(K_k\) and holds for any conditioning set.
\end{proof}

\Cref{cor:model_coherent_forecast} places the model with the methods that
obtain coherence by construction \citep{pennings2017integrated,
rangapuram2021hiere2e,das_topdown2023}, in contrast with the reconciliation
projections that map incoherent base forecasts onto the coherent subspace
\citep{hyndman2011optimal,wickramasuriya2019mint}. Two properties separate it
inside that family. The support is the integer lattice intersected with the
coherent set, so counts and durations stay whole and the zero atom is
preserved; and the conditional law at each resolution is indexed by a graph
\(\mathcal G^{(k)}\) that is learned, so the forecast carries a per-resolution
structural account alongside its predictive distribution.
